\section{Conclusion}
\label{sec:conclusion}

This paper presents a lightweight, training-free Text-to-SQL pipeline
that achieves EX\,=\,0.824 on Spider~1.0 using a 9B-parameter local LLM
and an INT8-quantized arctic-embed-m schema retriever.
Five pipeline components are each validated with statistically significant
improvements ($p < 0.001$, McNemar's test) in a controlled ablation study.

The central finding is that dynamic INT8 quantization of the schema
retrieval encoder introduces zero measurable degradation in retrieval
quality ($\Delta$R@5\,=\,0.0000, $\Delta$MRR\,=\,0.0000 on Spider~1.0;
$\Delta$R@10\,=\,0.0000 on Spider~2.0-Lite) while reducing encoder RAM
by 78.2\% (418\,MB $\to$ 91\,MB) and improving encoding latency by 15.2\%.
A 12-block per-layer sensitivity analysis confirms that uniform INT8
is the globally optimal quantization strategy for this model.

These results establish that INT8 quantization is a safe default for
dense retrieval encoders in Text-to-SQL pipelines, enabling
deployment on resource-constrained edge devices---Jetson Nano,
Raspberry Pi~4, mobile platforms---without retrieval-quality risk.

Validation on Spider~2.0-Lite confirms that the pipeline and
its quantization properties generalize to a substantially harder
benchmark, achieving a 55\% relative improvement in execution accuracy
through few-shot retrieval from a compact gold SQL pool.

Three directions appear most productive for future work.
Replacing the 9B LLM with a fine-tuned code model is expected to
substantially reduce semantic errors on analytical Spider~2.0-Lite queries.
Expanding the Spider~2.0-Lite few-shot pool beyond 24 examples, ideally
with per-database coverage, is expected to improve few-shot retrieval
effectiveness on this benchmark.
Finally, extending the INT8 sensitivity analysis to the LLM component
itself---using AWQ or GPTQ quantization of Qwen3.5---represents a natural
continuation of the quantization safety characterization initiated here.
